% ═════════════════════════════════════════════════════════════════════════════
\section{SMT-Based Methods}
\label{sec:smt}
% ═════════════════════════════════════════════════════════════════════════════

The rise of powerful Satisfiability Modulo Theories (SMT) solvers over the
past two decades has opened new avenues for parameterized verification.
SMT-based methods encode system states, transitions, and properties as logical
formulas and use highly optimized solvers to reason about them
automatically~\cite{BSST09}. Their main advantage over purely deductive
approaches is that the invariant search is largely automated; their main
advantage over finite-state model checking is that sets of states are
represented by formulas over rich theories, which makes them applicable to
systems with rich data and large process counts.

\subsection{Array-Based Systems}

In \emph{array-based systems}~\cite{GR10,GS10}, the parameterized state is
modeled as a collection of arrays indexed by process identifiers, and
quantified formulas express properties over all processes simultaneously. For
the running example, a single array~$s$ maps each process index to its local
state. The initial condition is the formula stating that for every index~$i$,
$s[i] = \Idle$. The transition from~$\Wait$ to~$\Crit$ is described by a
formula stating that there is an index~$k$ with $s[k] = \Wait$, that every
index $j \neq k$ satisfies $s[j] \neq \Crit$, and that the new array is
obtained from~$s$ by setting $s[k]$ to~$\Crit$. The bad configurations are
described by the formula stating that there are two distinct indices $i$
and~$j$ with $s[i] = \Crit$ and $s[j] = \Crit$.

The MCMT tool~\cite{GS10} implements model checking for array-based systems
via backward reachability in a quantified fragment of first-order logic, using
SMT solvers to discharge the fixpoint conditions. The sets of configurations
encountered during backward reachability are represented by existentially
quantified formulas, such as the formula for the bad configurations above. The
predecessors of such a formula are again existentially quantified, provided
that universal guards are handled appropriately; a common treatment
instantiates the universal quantifier only with the indices mentioned in the
formula, which over-approximates the system in a way that closely corresponds
to monotonic abstraction (Section~\ref{sec:rmc}). The test that a new formula
adds no configurations to those already found reduces to the satisfiability of
a formula with an $\exists^*\forall^*$ quantifier prefix, which SMT solvers
handle by instantiation.

In the running example, backward reachability starts from the formula for the
bad configurations. A step from~$\Wait$ to~$\Crit$ by some process~$k$ can
only lead to a bad configuration if the other process in~$\Crit$ was already
there before the step, which contradicts the guard. The other transitions
lead from a configuration to a bad one only if the configuration is itself
bad. The fixpoint is therefore reached immediately, and since no initial
configuration satisfies the formula, the system is safe.

This approach handles systems in which each process has a local state drawn
from an arbitrary (possibly infinite) domain, modeled as an interpreted theory
in the SMT solver. Termination of backward reachability is not guaranteed in
general, but it is guaranteed for classes of systems in which the models of
the existential formulas, viewed as finite structures, are well-quasi-ordered,
which connects the approach to the theory of Section~\ref{sec:wqo}~\cite{GR10}.
The Cubicle model checker~\cite{CGKMZ12} follows the same principles, with a
parallel implementation and a technique that uses finite instances of the
system to guess approximations of the reachable states, with backtracking
when an approximation turns out to be too coarse~\cite{CGKMZ13}. Array-based
methods have seen success on cache coherence and distributed protocols, and
have connections to the decidable array property fragments studied in program
verification~\cite{BMS06}.

\subsection{IC3/PDR and Quantified Invariants}

A prominent line of work in this space is \emph{IC3/PDR} (Property Directed
Reachability), originally developed for finite-state hardware
verification~\cite{Bra11,EMB11} and later extended to the parameterized
setting through quantified invariant inference. IC3/PDR works by incrementally
constructing an inductive strengthening of the safety property. It maintains
a sequence of formulas that over-approximate the sets of states reachable
within increasing numbers of steps. It strengthens these formulas by blocking
states from which a bad state can be reached, propagates facts from each
formula to the next, and stops when two consecutive formulas coincide, at
which point they form an inductive invariant.
In the parameterized setting, these formulas are universally quantified over
process identifiers, and SMT reasoning over the resulting quantified theories
replaces the finite propositional checks of the original algorithm.

Karbyshev et al.~\cite{KBIRS17} developed a variant of PDR that infers
universally quantified invariants. The \emph{diagram} of a state is an
existential formula that describes the state up to the identity of its
elements. When a state must be blocked, the algorithm blocks its diagram.
Blocking a diagram excludes the state together with all larger states that
contain it, which makes the method similar to backward reachability over
upward-closed sets. When it terminates, the algorithm returns a
counterexample, a universal inductive invariant, or a proof that no such
invariant exists.
Gurfinkel, Shoham, and Meshman~\cite{GSM16} exploit the symmetry of
parameterized systems to reduce the inference of quantified invariants to the
inference of invariants for a small number of representative processes, a task
that can be delegated to an SMT-based model checker.

\subsection{Decidable Logics and Ivy}

Tools such as Ivy~\cite{PMPSS16} build on these ideas to provide interactive and automated frameworks for specifying and verifying
distributed protocols. Ivy allows users to specify systems in a first-order
language and restricts the verification conditions to the \emph{effectively
propositional} fragment EPR, in which formulas have a quantifier prefix of the
form $\exists^*\forall^*$ and contain no function symbols, extended with
function symbols whose use satisfies an acyclicity condition. Satisfiability
in this fragment is decidable, and every satisfiable formula has
a finite model whose size is bounded by the formula. Consequently, every
failed check yields a small, finite counterexample, which the user can
inspect, as described in Section~\ref{sec:deductive}. The bound on the size
of models plays a role similar to a cutoff: it limits the number of processes,
messages, and values that need to be considered.

Many protocols do not fit into EPR directly. Padon et al.~\cite{PLSS17}
showed that complex consensus protocols, including several variants of Paxos,
can be verified within EPR by introducing derived relations that keep the
verification conditions in the decidable fragment. In practice, this workflow has successfully verified a range of
distributed protocols, including variants of Paxos, two-phase commit, and
leader election. The inference of invariants for such protocols has also been
automated. One approach generalizes invariants of small finite instances to
quantified invariants and checks them for all sizes~\cite{MGJKKS19}, in the
spirit of invisible invariants; another searches for quantified formulas that
separate reachable states from states that lead to an error~\cite{KPIA20}.

\subsection{Threshold-Based Distributed Algorithms}

Fault-tolerant distributed algorithms form a further class of parameterized
systems, in which the parameters include not only the number~$n$ of processes
but also an upper bound~$t$ on the number of faulty processes. Processes
typically wait until they have received messages from more than a threshold
of other processes, such as $(n+t)/2$, before taking a step. Such algorithms were first
verified by abstraction-based parameterized model checking~\cite{JKSVW13},
and have since been modeled as \emph{threshold automata}, in which
transitions are guarded by linear constraints over message counters and
parameters~\cite{KVW17}. A key result is a \emph{short counterexample
property}: if a threshold automaton violates a safety or liveness property for
some values of the parameters, then it does so in a run whose length, measured
in a suitable sense, is bounded independently of the parameters. Bounded model
checking with an SMT solver over the parameters then suffices for all
parameter values~\cite{KVW17,KLVW17}.

\subsection{Discussion}

The relationship between SMT-based methods and the other approaches in this
chapter is instructive. SMT methods can be seen as automating the invariant
discovery step that deductive methods (Section~\ref{sec:deductive}) leave to
the user, while using a logical representation of state that complements the
regular-language and upward-closed-set representations of
Section~\ref{sec:rmc}: logical formulas easily
express data and arithmetic, while automata easily express the ordering of
processes in a configuration. The tradeoff is decidability: the quantified
theories used by SMT-based tools are generally undecidable, and the tools rely
on heuristics, user interaction, or restriction to decidable fragments to
terminate. SMT-based methods are thus more automatic than purely deductive
approaches, but they lack the termination guarantees of cutoff and monotonic
abstraction methods.
